% ============================================================
% SECTION 5 — Benchmark Construction (Method)
% ============================================================
\section{Benchmark Construction}
\label{sec:method}

The benchmark is a controlled representation intervention on a fixed
flood substrate. The same folds, same solvers, same targets, and same
evaluation protocol are held constant across three representation
levels; only the feature set changes. This section describes what is
built (the ladder), how inference is disaggregated (encode then decode),
and how pre-training quality signals drive operational adaptation (the
gearbox).

% ------------------------------------------------------------------
\subsection{Encode then decode}
\label{sec:encode-decode}

Geospatial prediction is encode-then-decode. The \emph{encoder}
identifies and joins the geography: which polygon is which ZCTA, which
addresses fall in which boundary, which units are neighbors, which
crosswalk version is current. This is classification work---discrete,
combinatorial, and Fano-native---and it produces the feature substrate
that everything downstream consumes. The \emph{decoder} consumes that
substrate: regression to predict continuous targets (claims, presence,
high-water marks), then a decision layer that allocates resources under
operational constraints.

The RSCT certificate primitives split cleanly across the two halves:
\begin{itemize}
\item \textbf{Encoder side.} Signal purity $\alpha = R/(R+N)$ is
  grounded where Fano's inequality applies natively: boundary
  classification, crosswalk identification, and spatial join correctness
  are discrete channel problems. The geometric compatibility
  $\kappa_\text{geom}$ (Phase~0.5) is also encoder-side: it measures
  spatial connectivity, feature coverage, and scale stability from
  problem structure alone, before any model is trained.
\item \textbf{Decoder side.} The task residual floor ($\TRF$) estimator
  is the decoder-side analog of Fano's noise-floor role: a
  rate-distortion lower bound on achievable prediction error,
  regime-appropriate for continuous targets. Cross-validation fold
  metrics (R$^2$, RMSE, ROC-AUC) are decoder-side measurements.
\item \textbf{Decision layer.} The gate scaffolding---sequential
  gates producing \textsc{execute}, \textsc{caution}, or
  \textsc{refuse}---operates at the boundary between prediction and
  operational use, where the question shifts from ``how accurate?'' to
  ``accurate enough for this decision?''
\end{itemize}

This separation matters because the benchmark's central claim---that
audit flags predict which representation fixes help---requires
encoder-side signals ($\kappa_\text{geom}$, spatial structure) to
predict decoder-side outcomes (fold-level uplift). If both sides used
the same measurements, the prediction would be circular.

% ------------------------------------------------------------------
\subsection{The representation ladder}
\label{sec:ladder}

Three representation levels form a strict superset chain. No features
are removed or re-tuned between levels, so any performance change is
attributable to the added features alone.

\begin{description}
\item[R0 (static tabular).] ACS demographics, SVI vulnerability indices,
  FEMA flood zone percentages, terrain derivatives (TWI, slope), HIFLD
  infrastructure proximity, and temporally-gated NFIP claim history.
  Approximately 33--36 features depending on scenario.
\item[R1 (hydrology + spatial lag).] R0 plus NHDPlus catchment
  attributes, levee proximity, sewershed type (NYC), and eight W-matrix
  spatial lag features computed per fold from training data only to
  prevent target leakage (Appendix~\ref{app:relational-ablation}).
  Approximately 58--63 features.
\item[R2 (temporal event dynamics).] R1 plus MRMS precipitation timing,
  tide gauge anomaly, and HURDAT2 storm-track features. Approximately
  67--72 features. Temporal sources have vintage boundaries
  (Table~\ref{app:tab:availability}); events predating the source are
  marked \texttt{not\_applicable\_temporal}, not imputed.
\end{description}

Solvers are held constant across levels: HistGBDT
(\texttt{max\_depth=6}, \texttt{learning\_rate=0.1}, \texttt{max\_iter=200})
and Ridge ($\alpha{=}1.0$, median imputation, standard scaling). No
hyperparameter tuning is performed. Random seed is 42. Fold assignments
are generated once (spatial blocking by county or ZIP3 fallback) and
reused at all levels.

% ------------------------------------------------------------------
\subsection{The gearbox: pre-training quality signals}
\label{sec:gearbox}

The encode/decode separation creates a problem for operational routing:
$\kappa_\text{geom}$ is encoder-side and level-invariant---it measures
problem geometry, not model quality---so it cannot distinguish between
R0, R1, and R2 for a given cell. A routing system that uses
$\kappa_\text{geom}$ alone will always recommend the same arm regardless
of representation level.

The \emph{gearbox} (Phase~0.75) sits at the boundary between encode and
decode. It runs a cheap pre-training probe---Ridge cross-validation on
the R0 feature set---and derives five quality signals per
(scenario, target) cell:

\begin{description}
\item[$\tau$ (stability).] $\tau = 1/(1 + \text{CV})$ where
  $\text{CV} = \text{std}(\text{folds}) / |\text{mean}(\text{folds})|$.
  High $\tau$ means stable fold performance; low $\tau$ means the cell
  is sensitive to fold composition.
\item[$\sigma$ (turbulence).] $\sigma = \text{std}(\text{fold metrics},
  \text{ddof}{=}1)$. Raw fold variance, independent of mean level.
\item[$\alpha_\text{warmup}$ (signal strength).] For regression,
  $\alpha = \max(0, R^2)$; for classification,
  $\alpha = 2 \cdot (\text{AUC} - 0.5)$. Rescaled to $[0,1]$ so both
  task types share a common quality axis.
\item[Collapse risk.] Fraction of folds where the Ridge solver fails to
  beat the mean predictor ($R^2 \leq 0$ or AUC $\leq 0.5$).
\item[Coherence.] $1 - \text{IQR} / |\text{median}|$, measuring
  whether folds agree on the direction and scale of signal.
\end{description}

These five signals are then calibrated into a discrete \emph{gear
assignment} per cell, following the GearboxCalibrator methodology: tau
values are collected across all cells, winsorized, and partitioned into
percentile-based boundaries with monotonicity guards. Each cell receives
a gear from a five-level system (Precision, Balanced, Explore, Discovery,
Reverse) that determines how aggressively the decoder should operate.

The gearbox output feeds into Phase~6 (DGM routing), which evaluates
two strategies side by side:
\begin{enumerate}
\item \textbf{Kappa-first (legacy).} Uses $\kappa_\text{geom}$ as the
  primary discriminant. Known limitation: level-invariant, so it always
  exits on the first branch.
\item \textbf{Gear-informed (new).} Uses gearbox signals (gear,
  $\alpha_\text{warmup}$, $\sigma$) as the primary discriminant, falling
  back to certificate comparison only when signals are ambiguous. These
  signals vary per cell, enabling actual discrimination across arms.
\end{enumerate}

The gear-informed strategy produces internal decisions
(\textsc{execute}, \textsc{re\_encode}, \textsc{repair},
\textsc{fallback}) that are projected to public decisions
(\textsc{execute}, \textsc{caution}, \textsc{refuse}) at the output
boundary---two vocabularies deliberately separated so that system-internal
routing logic does not leak into user-facing operational guidance.

% ------------------------------------------------------------------
\subsection{Pre-registration and tamper evidence}
\label{sec:preregistration}

Diagnostics and $\kappa_\text{geom}$ are computed and uploaded to S3
with timestamps \emph{before} the next training level executes.
Phase~0.5 (geometry kappa) precedes Phase~1 (R0 training); Phase~0.75
(gearbox warmup) generates its own folds from spatial structure rather
than consuming folds produced by Phase~1, maintaining independence. S3
object timestamps provide tamper-evident ordering: the predictor
($\kappa_\text{geom}$, warmup signals) is frozen before the outcome
(R1/R2 uplift) exists.

The full pre-registered protocol---hypotheses, design matrix, controlled
variables, statistical tests, kill rules, and change control---is in
Appendix~\ref{app:doe}.
